% Options for packages loaded elsewhere
\PassOptionsToPackage{unicode}{hyperref}
\PassOptionsToPackage{hyphens}{url}
\PassOptionsToPackage{dvipsnames,svgnames,x11names}{xcolor}
%
\documentclass[
  11pt,
  a4paper,
]{ltjsarticle}
\usepackage{amsmath,amssymb}
\usepackage{iftex}
\ifPDFTeX
  \usepackage[T1]{fontenc}
  \usepackage[utf8]{inputenc}
  \usepackage{textcomp} % provide euro and other symbols
\else % if luatex or xetex
  \usepackage{unicode-math} % this also loads fontspec
  \defaultfontfeatures{Scale=MatchLowercase}
  \defaultfontfeatures[\rmfamily]{Ligatures=TeX,Scale=1}
\fi
\usepackage{lmodern}
\ifPDFTeX\else
  % xetex/luatex font selection
\fi
% Use upquote if available, for straight quotes in verbatim environments
\IfFileExists{upquote.sty}{\usepackage{upquote}}{}
\IfFileExists{microtype.sty}{% use microtype if available
  \usepackage[]{microtype}
  \UseMicrotypeSet[protrusion]{basicmath} % disable protrusion for tt fonts
}{}
\makeatletter
\@ifundefined{KOMAClassName}{% if non-KOMA class
  \IfFileExists{parskip.sty}{%
    \usepackage{parskip}
  }{% else
    \setlength{\parindent}{0pt}
    \setlength{\parskip}{6pt plus 2pt minus 1pt}}
}{% if KOMA class
  \KOMAoptions{parskip=half}}
\makeatother
\usepackage{xcolor}
\usepackage[margin=24mm]{geometry}
\setlength{\emergencystretch}{3em} % prevent overfull lines
\providecommand{\tightlist}{%
  \setlength{\itemsep}{0pt}\setlength{\parskip}{0pt}}
\setcounter{secnumdepth}{5}
\ifLuaTeX
  \usepackage{selnolig}  % disable illegal ligatures
\fi
\IfFileExists{bookmark.sty}{\usepackage{bookmark}}{\usepackage{hyperref}}
\IfFileExists{xurl.sty}{\usepackage{xurl}}{} % add URL line breaks if available
\urlstyle{same}
\hypersetup{
  colorlinks=true,
  linkcolor={blue},
  filecolor={Maroon},
  citecolor={Blue},
  urlcolor={blue},
  pdfcreator={LaTeX via pandoc}}

\author{}
\date{}

\begin{document}

\hypertarget{ux9023ux7acbux4e00ux6b21ux65b9ux7a0bux5f0f}{%
\section{連立一次方程式}\label{ux9023ux7acbux4e00ux6b21ux65b9ux7a0bux5f0f}}

\[
Ax=b
\]

を解く問題には3つの見方があります。

\begin{enumerate}
\def\labelenumi{\arabic{enumi}.}
\tightlist
\item
  方程式の集合として見る。
\item
  \texttt{A} の列の一次結合として見る。
\item
  線形写像 \texttt{A} の逆像を求める問題として見る。
\end{enumerate}

\hypertarget{ux89e3ux304cux5b58ux5728ux3059ux308bux6761ux4ef6}{%
\subsection{解が存在する条件}\label{ux89e3ux304cux5b58ux5728ux3059ux308bux6761ux4ef6}}

\texttt{A={[}a\_1\ ...\ a\_n{]}} とすると

\[
Ax=x_1a_1+\cdots+x_na_n
\]

なので、解が存在する条件は

\[
b\in\operatorname{Col}(A)
\]

です。

\hypertarget{ux89e3ux306eux5f62}{%
\subsection{解の形}\label{ux89e3ux306eux5f62}}

一つの特解 \texttt{x\_p} があれば、すべての解は

\[
x=x_p+z,\qquad z\in\ker A
\]

です。

つまり解集合は「一つの解 +
零空間」で表されます。この見方は自由変数の意味を明確にします。

\end{document}
